\section{Introduction}

A cryptographically relevant quantum computer would break RSA, finite-field
and elliptic-curve Diffie-Hellman, ECDSA and EdDSA through Shor's
algorithm~\cite{shor1997}. Traffic recorded today can be decrypted once such a
machine exists, so key exchange must move first~\cite{mosca2018}. NIST has
standardized ML-KEM, ML-DSA and SLH-DSA~\cite{fips203,fips204,fips205}, and
its transition report, NIST IR~8547~\cite{ir8547}, deprecates every
quantum-vulnerable public-key family after 2030 and disallows it after 2035.
In the United States the schedule is now an obligation: Executive
Order~14412~\cite{eo14412} sets federal migration deadlines, and OMB memorandum
M-26-15~\cite{omb2615} asks each agency for a migration plan aligned with
IR~8547.

Every such plan begins with an inventory. The cryptography an organization
runs is spread across layers that no single view captures: calls in source
code, key sizes and group lists in configuration files, certificates on disk,
libraries pulled in by lockfiles, code compiled into shipped binaries, and,
finally, what a server actually negotiates with a client. The open-source tools
available to build an inventory read one of these layers, almost always
source code. A source scanner that sees a post-quantum library import cannot
tell whether the deployed endpoint uses it.

We built \tool{} to close that gap and used it to measure where the migration
stands. \tool{} looks for cryptography in five layers and records everything it
finds in one format. It then links findings that belong to the same system and
reports where the layers disagree. The case it was designed for is the
\emph{declared-versus-deployed gap}: code or dependencies that support
post-quantum key exchange in front of an endpoint that still negotiates only
classical groups. Findings are checked against the IR~8547 and CNSA~2.0
schedules, ranked by a risk score, and exported as a CycloneDX cryptographic
bill of materials (CBOM)~\cite{cyclonedx} and as SARIF~\cite{sarif}.

We ask five questions.
\begin{itemize}
\item \RQ{1} (tool): Does scanning several layers and comparing them find more
  than a source-code scanner does, and which layers matter most?
\item \RQ{2} (federal web): What share of public federal \texttt{.gov} TLS
  endpoints negotiate hybrid post-quantum key exchange, and does that depend
  more on the agency or on the hosting provider?
\item \RQ{3} (open source): How common is quantum-vulnerable cryptography in
  the most critical open-source projects, and in which layers does it appear?
\item \RQ{4} (cost): What does hybrid key exchange add in bytes and
  computation?
\item \RQ{5} (declared vs.\ deployed): For federal services whose code is
  public, does what the code declares match what the deployed endpoint
  negotiates?
\end{itemize}

\paragraph{Contributions.}
\begin{itemize}
\item \tool{}, an open-source scanner covering five layers, with cross-layer
  reconciliation, policy evaluation and a GitHub Action, released under the
  Apache~2.0 license (\S\ref{sec:design}).
\item A labeled corpus of \data{rqone-cases} cases across five layers and
  \data{rqone-languages} language families, and a harness that runs
  \data{rqone-competitors} other open-source scanners in pinned containers.
  \tool{} finds \data{rqone-pqscope-file-r} of the labeled assets against
  \data{rqone-best-competitor-r} for the best alternative; restricted to source
  code the margin is smaller, and one competitor does better on Java
  (\S\ref{sec:rq1}).
\item A scan of all \data{rqtwo-endpoints} apex and \texttt{www} endpoints of
  the \data{rqtwo-domains} federal \texttt{.gov} domains, in which
  \data{rqtwo-ep-share} of TLS endpoints negotiate a hybrid group. It confirms
  for the federal population what studies of the wider web
  found~\cite{wickramasinghe2026,loizou2026}, that hosting drives adoption, and
  measures how much the agency adds once the provider is known
  (\S\ref{sec:rq2}).
\item A scan of the \data{rqthree-projects} most critical distinct open-source
  code bases by OpenSSF Criticality Score, in which \data{rqthree-any-share} of
  those scanned carry quantum-vulnerable cryptography, with an audit of a
  sample of findings that puts \tool{}'s precision on this code at \data{rqthree-audit-weighted}
  (\S\ref{sec:rq3}).
\item Benchmarks of the NIST post-quantum parameter sets against their
  classical counterparts, with every size cross-checked against OpenSSL
  (\S\ref{sec:rq4}).
\item A comparison of declared and deployed cryptography for
  \data{fed-scanned} public federal repositories and the endpoints they serve.
  The gap \tool{} was built to find never occurs; the reverse is common.
  \data{fed-hybrid-endpoints-share} of the endpoints negotiate a hybrid group,
  all of them through \data{fed-hybrid-providers} hosting providers, while
  \data{fed-nocrypto-share} of the repositories contain no public-key
  cryptography that \tool{} detects (\S\ref{sec:fed}).
\end{itemize}
Every number in the paper is generated by a script from the released data;
none is typed by hand.
